%
%	Szablon pracy dyplowowej inżynierskiej
%	rok akademicki 2020/21
%	Autor: Małgorzta Karabowicz
%	na podstawie szablonu stworzonego przez Przemysława Michalskiego (Wydział Fizyki PW)
%
%	brak gwarancji na 100% zgodność z najnowzym wzorem pracy inż. udostępnionym przez dziekanat
%


\documentclass[12pt,a4paper,twoside]{book}
\linespread{1.15} %interlinia

\usepackage{polski}

%czcionki
\usepackage[T1]{fontenc}
\usepackage[utf8]{inputenc}
 
\usepackage[scaled=0.92]{helvet} % użyta czcionka Helvetica
\renewcommand\familydefault{\sfdefault} %czcionka bezszeryfowa
\fontshape{n}
\selectfont

%definicje kolorów
\usepackage[table]{xcolor}
\definecolor{sliwka}{RGB}{150,95,119} %kolor œliwkowy, zgodny z zaleceniami
\definecolor{grafit}{RGB}{60,60,60} %kolor grafitowy, zgodny z zaleceniami (przy druku monochromatycznym)

%linki w PDFie - linkowane są przypisy oraz bibliografia
\usepackage{hyperref}
\hypersetup{
	pdftitle = {}, %tytuł dokumentu PDF
	pdfauthor = {}, %autor PDFa
	colorlinks,
	linktoc=page,
	linkcolor=sliwka,
	citecolor=sliwka,
	filecolor=sliwka,
	urlcolor=sliwka
}

\usepackage{amsmath,amsfonts,amssymb} %fonty matematyczne
\usepackage{icomma} %brak spacji po przecinku w trybie matematycznym (zwykle przecinek jest traktowany jako rozdział między kolejnymi trzema cyframi)
\usepackage{textcomp} %the package supports the Text Companion fonts, which provide many text symbols
\usepackage{gensymb} %stopien Celsjusza (\celcius), stopien (\degree), mikro (\micro), om (\textohm) i promil (\perthousand) w trybie tekstowym

\usepackage{graphicx}
\usepackage{placeins} %defines a \FloatBarrier, beyond which floats cannot pass. Przydatne, gdy obrazki/tabele nie wstawiają się tam, gdzie chcemy.
\usepackage[singlelinecheck=false, format=hang,justification=raggedright]{caption}

\newcommand{\source}[1]{\vspace{0pt}{\hfill \footnotesize{(\textit{{#1}})} } } % można dodać źródło obrazka \source{}


%marginesy
\usepackage[left=3cm,right=2cm,top=2.5cm,bottom=2.5cm]{geometry} 


\usepackage{fancyhdr}
\pagestyle{fancy}
\renewcommand{\headrulewidth}{0pt} %brak linii w nagłówku
\renewcommand{\footrulewidth}{1.5pt} %linia w stopce
\renewcommand{\footrule}{{\color{sliwka}%
	\vskip-\footruleskip\vskip-\footrulewidth
	\hrule width\headwidth height\footrulewidth\vskip\footruleskip}} %kolor linii w stopce

\fancyhead{} %pusty nagłówek
\cfoot{} %pusta śœrodkowa część stopki
\fancyfoot[LE,RO]{\thepage} %numer strony w stopce na parzystej stronie z lewej, na nieparzystej - z prawej

%listy - kształt punktów
\renewcommand\labelitemi{\color{sliwka}{\tiny$\blacksquare$}} %1. poziom
\renewcommand\labelitemii{\color{sliwka}{\tiny$\square$}} %2. poziom
\renewcommand\labelitemiii{\color{sliwka}{$\cdot$}} %3. poziom
%wcięcia w listach
\usepackage{enumitem}
\setlist[itemize]{noitemsep, topsep=0pt, leftmargin=1cm, itemindent=0cm}

%stylistyka
\widowpenalty=10000 %nie pozostawia wdów na końcu strony
\clubpenalty=10000 %nie pozostawia sierot
\brokenpenalty 10000 %nie dzieli stron, gdy przenoszony jest wyraz
\sloppy % zakaz wydłużania linii poza margines
\usepackage[none]{hyphenat} %bez przenoszenia wyrazów
\frenchspacing %,,kontynentalne" zwyczaje typograficzne

%odstêpy
\usepackage{setspace}
\setstretch{1.15} %odstęp między wierszami
\setlength{\parindent}{0pt} %szerokość wciecia akapitowego
\setlength{\parskip}{4pt} % odstęp za akapitem

%obrazki/tabele/równania
\usepackage{array}
\usepackage{multicol}
\usepackage{multirow}
\renewcommand*{\figurename}{Rysunek} 
\renewcommand*{\tablename}{Tabela}
\setlength{\arrayrulewidth}{1pt} %gruboœść linii tabeli
\arrayrulecolor{sliwka} %kolor linii tabeli
\renewcommand\thefootnote{\textcolor{sliwka}{\arabic{footnote}}} %kolor numeru przypisu
\usepackage[font={footnotesize}]{caption} %czcionka 9 pkt w podpisach tabeli/obrazków
\numberwithin{equation}{chapter} %numerowanie równań od numeru rozdziału
\newcolumntype{M}[1]{>{\centering\arraybackslash}m{#1}}
\usepackage{subfig}



\let\lll=\relax


\graphicspath{ {./img/} {./rysunki/} {./wykresy/} }


% ustawienia tytułów rozdziałów
\usepackage{titlesec}
\usepackage{titlesec}
\titleformat{\chapter}[display]
  {\normalfont\bfseries}{}{0pt}{\Huge}
\titlespacing{\section}{0pt}{15pt}{10pt}
\titlespacing{\subsection}{0pt}{15pt}{10pt}



\begin{document}

\pagestyle{empty}

\begin{titlepage}
	
\begin{minipage}{\textwidth}
	\begin{center}
		\includegraphics[width=\textwidth]{logo.png} \\
			\vspace{3cm}
		\includegraphics[width=\textwidth]{inzynierska.png} \\ %podmieniæ na magisterska.png dla studiów II stopnia
			\vspace{1.5cm}
  			na kierunku fizyka techniczna \\
			w specjalności (specjalność) \\
			\vspace{2cm}	
		{\Large
			Tytuł pracy}

			\vspace{2cm}	
		{\huge
			Imię i nazwisko} \\
			Numer albumu: xxxxxx \\
			\vspace{1.5cm}	
			promotor: \\
			tytuł naukowy, imię i nazwisko \\
			\vspace{3cm}
			WARSZAWA 202x		
	\end{center}
\end{minipage}

\end{titlepage}

\newpage
\mbox{ }

\newpage
{\large \textbf{Streszczenie}}\\
{\color{sliwka}\rule[1pt]{\textwidth}{1.5pt}}

\begin{description}[leftmargin=2.2cm,font=\normalfont]
\item[Tytuł pracy:]  Tytuł
\end{description}

Streszczenie



\begin{flushright}
	Imię i nazwisko
\end{flushright}

\vspace{0.025\textheight}
\textit{sloowa kluczowe:}\\
\textit{szablon, LaTeX, praca dyplomowa}


\vspace{2cm}

\begin{multicols}{2}
	\begin{flushleft}
		(podpis opiekuna naukowego)
	\end{flushleft}
	\begin{flushright}
		(podpis dyplomanta)
	\end{flushright}
\end{multicols}

\newpage
\mbox{ }
\newpage
{\large \textbf{Abstract}}\\
{\color{sliwka}\rule[1pt]{\textwidth}{1.5pt}}


\begin{description}[leftmargin=3.2cm,font=\normalfont]
\item[Title of the thesis:] Title
\end{description}

Abstract



\vspace{0.025\textheight}
\textit{Keywords:}\\
\textit{template, LaTeX, thesis}


\newpage
\mbox{ }

\newpage
{\large \textbf{Oświadczenie o samodzielności wykonania pracy}}\\
{\color{sliwka}\rule[1pt]{\textwidth}{1.5pt}}

\vspace{0.6cm}

\begin{figure}[h!]
	\includegraphics[width=0.5\textwidth]{img/pw_logo_tekst.png}
\end{figure}

\vspace{0.6cm}
{\large imię i nazwisko}\\
nr indeksu\\
kierunek studiów\\

\begin{center}
	{\large \textbf{Oświadczenie}}
\end{center}

{\large Świadomy/-a odpowiedzialności karnej za składanie fałszywych zeznań oświadczam, że niniejsza praca dyplomowa została napisana przeze mnie samodzielnie, pod opieką kierującego pracą dyplomową.}


Jednocześnie oświadczam, że:
\begin{itemize}
	\item niniejsza praca dyplomowa nie narusza praw autorskich w rozumieniu ustawy z dnia 4 lutego 1994 roku o prawie autorskim i prawach pokrewnych (Dz.U. z 2006 r. Nr 90, poz. 631 z późn. zm.) oraz dóbr osobistych chronionych prawem cywilnym,
	\item niniejsza praca dyplomowa nie zawiera danych i informacji, które uzyskałem/-am w sposób niedozwolony,
	\item niniejsza praca dyplomowa nie była wcześniej podstawą żadnej innej urzędowej procedury związanej z nadawaniem dyplomów lub tytułów zawodowych,
	\item wszystkie informacje umieszczone w niniejszej pracy, uzyskane ze źródeł pisanych i elektronicznych, zostały udokumentowane w wykazie literatury odpowiednimi odnośnikami,
	\item znam regulacje prawne Politechniki Warszawskiej w sprawie zarządzania prawami autorskimi i prawami pokrewnymi, prawami własności przemysłowej oraz zasadami komercjalizacji.
\end{itemize}

\vspace{3cm}

\begin{multicols}{2}
	\begin{flushleft}
		Warszawa, dnia (data)
	\end{flushleft}
	\begin{flushright}
		(czytelny podpis dyplomanta)
	\end{flushright}
\end{multicols}

\newpage
\mbox{ }

\newpage
{\large \textbf{Oświadczenie o udzieleniu Uczelni licencji do pracy}}\\
{\color{sliwka}\rule[1pt]{\textwidth}{1.5pt}}

\vspace{0.6cm}

\begin{figure}[h]
	\includegraphics[width=0.5\textwidth]{img/pw_logo_tekst.png}
\end{figure}

\vspace{0.6cm}
{\large imię i nazwisko}\\
nr indeksu\\
kierunek studiów\\

\begin{center}
	{\large \textbf{Oświadczenie studenta w przedmiocie udzielenia licencji Politechnice Warszawskiej}}
\end{center}


Oświadczam, że jako autor / współautor* pracy dyplomowej pt.
\begin{center}
\textit{Tytuł pracy}
\end{center}
udzielam / nie udzielam* Politechnice Warszawskiej nieodpłatnej licencji na niewyłączne, nieograniczone w czasie, umieszczenie pracy dyplomowej w elektronicznych bazach danych oraz udostępnianie pracy dyplomowej w zamkniętym systemie bibliotecznym Politechniki Warszawskiej osobom zainteresowanym.
Licencja na udostępnienie pracy dyplomowej nie obejmuje wyrażenia zgody na wykorzystywanie pracy dyplomowej na żadnym innym polu eksploatacji, w szczególności kopiowania pracy dyplomowej w całości lub w części, utrwalania w innej formie czy zwielokrotniania.

\vspace{3cm}

\begin{multicols}{2}
	\begin{flushleft}
		Warszawa, dnia (data)
	\end{flushleft}
	\begin{flushright}
		(czytelny podpis dyplomanta)
	\end{flushright}
\end{multicols}
* - niepotrzebne skreślić

\newpage
\mbox{ }



%spis treœci
\tableofcontents
\addtocontents{toc}{\protect\thispagestyle{empty}}
\pagestyle{empty}

\chapter{Wstęp}
\thispagestyle{fancy}
%
% WSTĘP
%


\section{Section 1}
	Tekst...
\subsection{Subsection 1}


	\begin{figure}[th!]		
		\includegraphics[width=1\textwidth]{logonew.png}		
		\source{tutaj podaję źródło}
		\caption{Podpis obrazka}
        \label{particlesSM}
	\end{figure}
	%\FloatBarrier % kiedy odkomentujemy obrazek wstawi się dokładnie w tym miejscu, w którym chcemy
%\pagestyle{fancy} % można włączyć jak chcemy, aby puste strony również miały numer i linię w dolnej stopce

\chapter{Kolejny rozdział}
\thispagestyle{fancy}
Nowy rozdział



\begin{thebibliography}{99}
	% Zmieniæ na 99, w przypadku, gdy bibliografia liczy wiêcej ni¿ 9 pozycji.
	% W przypadku du¿ej liczby pozycji, wygodniej korzystaæ z BibTeXa.
	% Zalecany wzór wpisu bibliograficznego dla publikacji z czasopisma: nazwiska i inicja³y imion (przy dwóch imionach bez spacji po kropce!) autorów pisane s¹ kapitalikami. Nastêpnie dwukropek i tytu³ pracy pisany kursyw¹. Nastêpnie kropka i nazwa czasopisma, volume (pogrubiony), rok i numery stron. W przypadku podawania przedzia³ów liczbowych u¿ywa siê pó³pauzy (--) bez spacji, a nie ³¹cznika (-)!
	% Zalecany wzór wpisu dla ksi¹¿ki: nazwiska i inicja³y imion kapitalikami; nastêpnie dwukropek i tytu³ kursyw¹. Po kropce numer wydania (opcjonalnie), nazwa wydawnictwa, miejsce i rok wydania.
	\thispagestyle{fancy}
	\bibitem{yamada} \textsc{Yamada A., Iwane N., Harada Y., Nishimura S., Koyama Y., Tanaka I.}: \textit{Lithium Iron Borates as High-Capacity Battery Electrodes}. Advanced Materials \textbf{22} (2010) 3583--3587.
	\bibitem{pusheen} \texttt{http://www.pusheen.com/}, dostêp 1.02.2016 r.
	\bibitem{chemia-nieorg} \textsc{Durrant P.J., Durrant B.}: \textit{Zarys wspó³czesnej chemii nieorganicznej}. Wydanie pierwsze. Pañstwowe Wydawnictwa Naukowe, Warszawa 1965.
\end{thebibliography}


\listoffigures
\thispagestyle{fancy}

\listoftables
\thispagestyle{fancy}


\end{document}














